\section{Phương pháp và thuật toán đề xuất}
\label{sec:method}

Chương này trình bày hai mô hình đề xuất. Graph-Regularized Neighborhood Regression (GRNR) thuộc nhóm giải trực tiếp: mỗi sản phẩm được mô tả bằng một bài toán hồi quy có nghiệm đóng, không có vòng lặp huấn luyện. Diffusion-Filtered History Factorization (DFHF) thuộc nhóm cần huấn luyện: các bảng embedding sản phẩm được học bằng gradient. Hai mô hình được thiết kế và đánh giá độc lập; mỗi mô hình chỉ được so sánh với các phương pháp cùng nhóm ở Chương~\ref{sec:experiments}. Mọi công thức trong chương được viết theo đúng mã nguồn và cấu hình đã khóa của dự án.

\subsection{Mô hình bài toán}
\label{sec:problem}

Gọi $\mathcal{U}$ là tập người dùng và $\mathcal{I}$ là tập sản phẩm. Dữ liệu huấn luyện là ma trận tương tác nhị phân $X \in \{0,1\}^{|\mathcal{U}|\times|\mathcal{I}|}$, trong đó $X_{ui}=1$ nếu người dùng $u$ đã đánh giá tích cực sản phẩm $i$. Lịch sử của người dùng $u$ là $H_u=\{i : X_{ui}=1\}$ với độ dài $n_u=|H_u|$; số tương tác của sản phẩm $i$ là $c_i=\sum_u X_{ui}$. Với mỗi người dùng, mô hình gán điểm $s(u,i)$ cho các sản phẩm chưa thấy và trả về $K$ sản phẩm có điểm cao nhất.

Khi đánh giá, cần phân biệt bốn đối tượng. \emph{Tập huấn luyện gốc} là toàn bộ tương tác train và chỉ dùng để loại các sản phẩm người dùng đã tương tác khỏi danh sách gợi ý. \emph{Tập huấn luyện giữ lại} ở tỉ lệ $d\in\{1;0{,}75;0{,}5;0{,}25\}$ là đầu vào thực sự của mô hình. \emph{Tập ứng viên} gồm các sản phẩm có ít nhất một tương tác trong tập giữ lại. \emph{Nhóm người dùng warm} gồm những người dùng có sản phẩm mục tiêu trong validation thuộc tập ứng viên. Tỉ lệ $d$ là phần cạnh được giữ lại, khác với mật độ tuyệt đối $|E|/(|\mathcal{U}||\mathcal{I}|)$ của ma trận.

Điểm chung của hai mô hình là không có vector riêng cho từng người dùng: người dùng được biểu diễn hoàn toàn bằng lịch sử $H_u$. Vì vậy, khi có một lịch sử mới, mô hình chấm điểm được ngay mà không cần huấn luyện lại. Ngoài điểm chung này, hai mô hình khác nhau về cách học tham số và thuộc hai nhóm phương pháp khác nhau (\tabref{tab:groups}).

\begin{table}[htbp]
\centering
\caption[Hai nhóm phương pháp và mô hình đề xuất tương ứng]{Hai nhóm phương pháp và mô hình đề xuất tương ứng.}
\label{tab:groups}
\begin{tabularx}{\textwidth}{@{}lll>{\raggedright\arraybackslash}X@{}}
\toprule
Nhóm & Mô hình đề xuất & Baseline cùng nhóm & Cách học tham số \\
\midrule
Giải trực tiếp & GRNR & EASE, GF-CF & Nghiệm đóng hoặc lọc tuyến tính, không có epoch \\
Cần huấn luyện & DFHF & LightGCN, XSimGCL, DirectAU & Tối ưu bằng gradient trong 200 epoch \\
\bottomrule
\end{tabularx}
\end{table}

\subsection{Graph-Regularized Neighborhood Regression}
\label{sec:grnr}

Graph-Regularized Neighborhood Regression (GRNR) dự đoán khả năng người dùng thích sản phẩm $i$ bằng một tổ hợp tuyến tính của các sản phẩm láng giềng của $i$ có trong lịch sử người dùng. Mỗi sản phẩm có một bài toán hồi quy riêng, nhỏ, giải được chính xác. Ý tưởng hồi quy giữa các sản phẩm kế thừa từ SLIM~\cite{ning2011slim} và EASE~\cite{steck2019ease}; khác biệt của GRNR nằm ở việc giới hạn láng giềng, trọng số cho dữ liệu ngầm và các thành phần điều chuẩn theo đồ thị được mô tả dưới đây. \figref{fig:grnr-diagram} tóm tắt hai giai đoạn của mô hình.

\begin{figure}[htbp]
\centering
\begin{tikzpicture}[node distance=1.25cm and 0.42cm, every node/.style={font=\footnotesize}, box/.style={text width=2.75cm, minimum height=1.25cm}]
\node[chthreeblue, box] (x) {Ma trận tương tác};
\node[chthreeblue, box, right=of x] (nb) {Chọn tối đa 128 láng giềng theo lift};
\node[chthreeblue, box, right=of nb] (graph) {Đồ thị sản phẩm và Laplacian};
\node[chthreeblue, box, right=of graph] (solve) {Giải hồi quy cho từng sản phẩm (Cholesky)};
\node[chthreegreen, box, below=of x] (hist) {Lịch sử người dùng};
\node[chthreegreen, box, right=of hist] (score) {Chấm điểm bằng hệ số láng giềng và chặn};
\node[chthreegreen, box, right=of score] (tilt) {Giảm nhẹ sản phẩm phổ biến, loại sản phẩm đã xem};
\node[chthreeorange, box, right=of tilt] (topk) {Top-$K$ kèm phân rã đóng góp};
\draw[chthreeflow] (x) -- (nb);
\draw[chthreeflow] (nb) -- (graph);
\draw[chthreeflow] (graph) -- (solve);
\draw[chthreedash] (solve.south) -- ++(0,-0.38) -| node[pos=0.25, above, font=\scriptsize] {hệ số $\beta$, chặn $b$} (score.north);
\draw[chthreeflow] (hist) -- (score);
\draw[chthreeflow] (score) -- (tilt);
\draw[chthreeflow] (tilt) -- (topk);
\end{tikzpicture}
\caption[Sơ đồ hoạt động của Graph-Regularized Neighborhood Regression]{Sơ đồ hoạt động của Graph-Regularized Neighborhood Regression. Hàng trên (xanh dương) chạy một lần khi fit; hàng dưới (xanh lá) chạy cho mỗi người dùng, kể cả người dùng có lịch sử mới.}
\label{fig:grnr-diagram}
\end{figure}


\subsubsection{Lựa chọn láng giềng và đồ thị sản phẩm}
\label{sec:grnr-neighbors}

Gọi $C=X^\top X$ là ma trận đồng xuất hiện (đường chéo đặt bằng 0). Với sản phẩm đích $i$, đặt xác suất nền đã làm trơn
\begin{equation}
p_i = \frac{c_i+1}{|\mathcal{U}|+2}.
\end{equation}
Với mỗi sản phẩm nguồn $j$, mức đồng xuất hiện \emph{vượt kỳ vọng} và độ nâng (lift) được tính bởi
\begin{equation}
e_{ij} = C_{ji} - c_j\,p_i, \qquad \mathrm{lift}_{ij} = \frac{e_{ij}}{c_j+\alpha}, \quad \alpha = 10.
\label{eq:grnr-lift}
\end{equation}
Tập láng giềng $N_i$ gồm tối đa $k=128$ nguồn có $\mathrm{lift}_{ij}>0$, xếp theo $e_{ij}$ giảm dần (hòa thì theo chỉ số sản phẩm). Như vậy chỉ những sản phẩm xuất hiện cùng $i$ nhiều hơn mức ngẫu nhiên mới được dùng làm đặc trưng.

Từ các trọng số lift dương, xây dựng đồ thị sản phẩm có hướng $W^{\mathrm{dir}}_{ij}=\mathrm{lift}_{ij}$ với $j\in N_i$, đối xứng hóa bằng $W=\max(W^{\mathrm{dir}}, W^{\mathrm{dir}\top})$ và lấy Laplacian chuẩn hóa
\begin{equation}
L = I - D_W^{-1/2}\, W\, D_W^{-1/2},
\end{equation}
trong đó $D_W$ là ma trận bậc của $W$. Laplacian này được dùng làm điều chuẩn ở mục tiếp theo, để các sản phẩm láng giềng có liên hệ mạnh với nhau nhận hệ số gần nhau.

\subsubsection{Hàm mục tiêu và phương pháp giải}
\label{sec:grnr-objective}

Với sản phẩm $i$, mỗi người dùng $u$ được mô tả bởi vector đặc trưng gồm các láng giềng có trong lịch sử, chuẩn hóa theo độ dài lịch sử, và một hệ số chặn:
\begin{equation}
\phi_u^{(i)} = \left[\, \frac{X_{u,N_i}}{\sqrt{\max(n_u - X_{ui},\,1)}} \;;\; 1 \,\right] \in \mathbb{R}^{|N_i|+1}.
\label{eq:grnr-feature}
\end{equation}
Mẫu số dùng $n_u-X_{ui}$: khi người dùng đã tương tác với $i$, chính $i$ được loại khỏi độ dài lịch sử, để đặc trưng giống với lúc dự đoán một sản phẩm chưa thấy. Tham số $\theta_i=[\beta_i; b_i]$ gồm hệ số láng giềng $\beta_i$ và hệ số chặn $b_i$, được ước lượng bằng bình phương tối thiểu có trọng số trên \emph{toàn bộ} người dùng, theo tinh thần của mô hình phản hồi ngầm~\cite{hu2008wmf}:
\begin{equation}
\min_{\theta_i}\; \frac12\sum_{u\in\mathcal{U}} w_{ui}\bigl(X_{ui} - \phi_u^{(i)\top}\theta_i\bigr)^2 \;+\; \frac12(\theta_i-\theta_i^0)^\top M P_i\,(\theta_i-\theta_i^0),
\label{eq:grnr-objective}
\end{equation}
với $w_{ui}=10$ nếu $X_{ui}=1$ và $w_{ui}=1$ trong trường hợp ngược lại, $M=|\mathcal{U}|$. Ma trận phạt có dạng khối
\begin{equation}
P_i = \begin{bmatrix} \eta\,\mathrm{diag}(s_j)_{j\in N_i} + \lambda_g\bigl(L_{N_iN_i} - (1-\beta)B_i\bigr) & 0 \\ 0 & \eta_b \end{bmatrix},
\label{eq:grnr-precision}
\end{equation}
trong đó $\eta=10^{-5}$, $\eta_b=10^{-3}$, $\lambda_g=10^{-3}$, $\beta=0{,}5$. Hệ số $s_j=\min\bigl((c_j/(c_j+2))^{-2},\,10\bigr)$ phạt nặng hơn các láng giềng có ít dữ liệu. $L_{N_iN_i}$ là khối con của Laplacian trên các láng giềng, còn $B_i$ là ma trận đường chéo đo phần năng lượng ``biên'' của mỗi láng giềng, tức phần liên kết đồ thị trỏ ra ngoài $N_i$; hệ số $\beta$ giữ lại một nửa phần biên này. Điểm neo $\theta_i^0$ có hệ số láng giềng bằng 0 và hệ số chặn bằng nghiệm của bài toán chỉ có hằng số:
\begin{equation}
b_i^0 = \frac{w\,p_i}{1+(w-1)\,p_i}, \quad w=10.
\end{equation}

Hàm \eqref{eq:grnr-objective} là toàn phương và lồi chặt khi $P_i$ xác định dương, nên nghiệm duy nhất thỏa phương trình chuẩn
\begin{equation}
\bigl(\Phi_i^\top W_i \Phi_i + M P_i\bigr)\,\theta_i = \Phi_i^\top W_i\, x_i + M P_i\,\theta_i^0,
\label{eq:grnr-normal}
\end{equation}
với $\Phi_i$ là ma trận xếp các hàng $\phi_u^{(i)}$, $W_i=\mathrm{diag}(w_{ui})$ và $x_i$ là cột $i$ của $X$. Hệ \eqref{eq:grnr-normal} có kích thước $(|N_i|+1)\times(|N_i|+1)\le 129\times129$ và được giải bằng phân rã Cholesky ở độ chính xác float64. Cài đặt kiểm tra sai số dư tương đối $\|A\theta_i-r\|/\|r\|<10^{-8}$ cho từng sản phẩm. Sản phẩm không có tương tác nào trong tập giữ lại nhận nghiệm $\theta_i=\theta_i^0$.

\subsubsection{Thuật toán huấn luyện và suy luận}
\label{sec:grnr-algorithm}

Ma trận Gram $\Phi_i^\top W_i\Phi_i$ được tính hiệu quả nhờ nhận xét: phần lớn người dùng có $X_{ui}=0$ và trọng số 1, nên có thể dùng chung ma trận Gram toàn cục $G=\tilde X^\top\tilde X$ của lịch sử chuẩn hóa $\tilde X_{uj}=X_{uj}/\sqrt{\max(n_u,1)}$, rồi chỉ thay chính xác các hàng của người dùng dương tính (đặc trưng loại $i$ khỏi mẫu số, trọng số 10). Thuật toán~\ref{alg:grnr} tóm tắt quá trình.

\begin{algorithm}[htbp]
\caption[Huấn luyện và chấm điểm Graph-Regularized Neighborhood Regression]{Huấn luyện và chấm điểm Graph-Regularized Neighborhood Regression.}
\label{alg:grnr}
\KwIn{Ma trận tương tác giữ lại $X$; siêu tham số ở \tabref{tab:grnr-hparams}}
\KwOut{Ma trận hệ số thưa $\mathcal{B}\in\mathbb{R}^{|\mathcal{I}|\times|\mathcal{I}|}$, vector chặn $b$}
Tính $c$, $p$, $C=X^\top X$, lịch sử chuẩn hóa $\tilde X$ và Gram toàn cục $G=\tilde X^\top \tilde X$\;
\ForEach{sản phẩm $i$}{
  Chọn $N_i$ theo \eqref{eq:grnr-lift}; ghi trọng số lift vào đồ thị $W$\;
}
Tính Laplacian chuẩn hóa $L$ của đồ thị $W$\;
\ForEach{sản phẩm $i$}{
  \eIf{$c_i=0$}{$\theta_i\leftarrow\theta_i^0$\;}{
  Lấy khối Gram của $N_i$ từ $G$, thay chính xác các hàng của người dùng có $X_{ui}=1$\;
  Lập $P_i$ theo \eqref{eq:grnr-precision}; giải \eqref{eq:grnr-normal} bằng Cholesky; kiểm tra sai số dư\;}
  Ghi $\beta_i$ vào cột $i$ của $\mathcal{B}$ và $b_i$ vào $b$\;
}
\textbf{Chấm điểm:} $s(u,i)=\tilde X_{u,:}\,\mathcal{B}_{:,i}+b_i$; khi xếp hạng dùng $s(u,i)\cdot(1+c_i)^{-0{,}1}$\;
\end{algorithm}

Khi suy luận, điểm của người dùng $u$ cho sản phẩm $i$ là
\begin{equation}
s(u,i)=\sum_{j\in H_u\cap N_i}\beta_{ij}/\sqrt{n_u}+b_i.
\end{equation}
Riêng ở bước xếp hạng, điểm được nhân với hệ số $(1+c_i)^{-0{,}1}$ để giảm nhẹ ưu thế của sản phẩm phổ biến; hệ số này không tham gia vào bài toán \eqref{eq:grnr-objective}. Phép nhân chỉ giữ đúng chiều khi điểm dương. Kiểm tra trên 3.000 người dùng ngẫu nhiên ở $d=1$ cho thấy không có điểm âm nào trong top-20, nên trong thực nghiệm chiều tác động của hệ số này không bị đảo.

Vì điểm là tổng tuyến tính, có thể phân rã $s(u,i)$ thành đóng góp $\beta_{ij}/\sqrt{n_u}$ của từng sản phẩm $j$ trong lịch sử. Hệ thống demo (Mục~\ref{sec:demo}) dùng phân rã này để giải thích gợi ý. Đây là phân rã đại số của điểm số, không phải giải thích nhân quả về lý do người dùng thích sản phẩm.

\begin{table}[htbp]
\centering
\caption[Siêu tham số của Graph-Regularized Neighborhood Regression]{Siêu tham số của Graph-Regularized Neighborhood Regression. Cấu hình đã khóa của mô hình.}
\label{tab:grnr-hparams}
\begin{tabular}{@{}llr@{}}
\toprule
Ký hiệu & Ý nghĩa & Giá trị \\
\midrule
$k$ & Số láng giềng tối đa & 128 \\
$\alpha$ & Làm trơn mẫu số của lift & 10 \\
$w$ & Trọng số mẫu dương & 10 \\
$\eta$ & Phạt ridge cho hệ số láng giềng & $10^{-5}$ \\
$\alpha_s$, giới hạn & Làm trơn và giới hạn trên của hệ số $s_j$ & 2; 10 \\
$\lambda_g$ & Trọng số điều chuẩn đồ thị & $10^{-3}$ \\
$\beta$ & Phần năng lượng biên được giữ & 0,5 \\
$\eta_b$ & Phạt cho hệ số chặn & $10^{-3}$ \\
$\gamma$ & Số mũ của hệ số độ phổ biến khi xếp hạng & $-0{,}1$ \\
\bottomrule
\end{tabular}
\end{table}

\subsection{Diffusion-Filtered History Factorization}
\label{sec:dfhf}

Diffusion-Filtered History Factorization (DFHF) biểu diễn người dùng bằng trung bình các embedding sản phẩm trong lịch sử, tương tự ý tưởng của FISM~\cite{kabbur2013fism} và mô hình bất đối xứng của Koren~\cite{koren2008factorization}. Khác biệt của DFHF là các bảng embedding được làm trơn bằng một bộ lọc khuếch tán trên đồ thị người dùng--sản phẩm trước khi dùng, và mô hình được huấn luyện bằng softmax trên toàn bộ danh mục. \figref{fig:dfhf-diagram} tóm tắt quá trình huấn luyện và gợi ý.

\begin{figure}[htbp]
\centering
\begin{tikzpicture}[node distance=1.25cm and 0.42cm, every node/.style={font=\footnotesize}, box/.style={text width=2.75cm, minimum height=1.25cm}]
\node[chthreeblue, box] (sv) {Hai bảng embedding: nguồn $S$, đích $V$};
\node[chthreeblue, box, right=of sv] (filt) {Bộ lọc khuếch tán trên đồ thị};
\node[chthreeblue, box, right=of filt] (ctx) {Ngữ cảnh: lịch sử bỏ sản phẩm cần dự đoán};
\node[chthreeblue, box, right=of ctx] (loss) {Softmax trên toàn danh mục};
\node[chthreegreen, box, below=of sv] (hist) {Lịch sử người dùng};
\node[chthreegreen, box, right=of hist] (pool) {Trung bình embedding nguồn đã lọc};
\node[chthreegreen, box, right=of pool] (score) {Tích với embedding đích, cộng chặn};
\node[chthreeorange, box, right=of score] (topk) {Top-$K$, loại sản phẩm đã xem};
\draw[chthreeflow] (sv) -- (filt);
\draw[chthreeflow] (filt) -- (ctx);
\draw[chthreeflow] (ctx) -- (loss);
\draw[chthreedash] (loss.north) -- ++(0,0.32) -| node[pos=0.25, above, font=\scriptsize] {cập nhật Adam, 200 epoch} (sv.north);
\draw[chthreedash] (filt.south) -- node[right, font=\scriptsize] {$\tilde S,\tilde V,b$} (pool.north);
\draw[chthreeflow] (hist) -- (pool);
\draw[chthreeflow] (pool) -- (score);
\draw[chthreeflow] (score) -- (topk);
\end{tikzpicture}
\caption[Sơ đồ hoạt động của Diffusion-Filtered History Factorization]{Sơ đồ hoạt động của Diffusion-Filtered History Factorization. Hàng trên (xanh dương) là vòng huấn luyện; hàng dưới (xanh lá) là bước gợi ý, dùng toàn bộ lịch sử. Người dùng chỉ có một tương tác nhận vector lịch sử bằng 0 và được xếp hạng bằng chặn $b$.}
\label{fig:dfhf-diagram}
\end{figure}


\subsubsection{Biểu diễn lịch sử và bộ lọc khuếch tán}
\label{sec:dfhf-representation}

DFHF có hai bảng embedding sản phẩm không ràng buộc với nhau, mỗi bảng có $64$ chiều: bảng nguồn $S\in\mathbb{R}^{|\mathcal{I}|\times 64}$ dùng để mô tả lịch sử và bảng đích $V\in\mathbb{R}^{|\mathcal{I}|\times 64}$ dùng để mô tả sản phẩm được chấm điểm, cùng một vector chặn $b\in\mathbb{R}^{|\mathcal{I}|}$. Gọi $D_U$, $D_I$ là ma trận bậc người dùng và sản phẩm của $X$. Toán tử chuyển
\begin{equation}
T = D_I^{-1}X^\top D_U^{-1}X
\end{equation}
lấy trung bình một vector trên các người dùng đã tương tác với sản phẩm, sau đó trên các sản phẩm của những người dùng đó; $T$ là ma trận ngẫu nhiên theo hàng của một bước đi ngẫu nhiên hai bước sản phẩm $\to$ người dùng $\to$ sản phẩm. Bộ lọc đa thức
\begin{equation}
F = \tfrac12 I + \tfrac14 T + \tfrac14 T^2
\label{eq:dfhf-filter}
\end{equation}
được áp dụng cho cả hai bảng: $\tilde S=FS$, $\tilde V=FV$. Sản phẩm không có tương tác giữ nguyên vector gốc. Biểu diễn người dùng và điểm số là
\begin{equation}
q_u = \frac{1}{n_u^{1/2}}\sum_{j\in H_u}\tilde S_j, \qquad s(u,i)=q_u^\top\tilde V_i + b_i.
\label{eq:dfhf-score}
\end{equation}
Theo cấu hình đã khóa, người dùng có đúng một tương tác trong tập giữ lại nhận $q_u=0$ (trọng số lịch sử đơn bằng 0), cả khi huấn luyện lẫn khi suy luận. Với những người dùng này, DFHF xếp hạng chỉ bằng $b_i$, tức một thứ tự độ phổ biến đã học. Hệ quả này rất quan trọng khi tập huấn luyện bị giảm mạnh: ở $d=0{,}25$, 68\% người dùng warm chỉ còn một tương tác (Mục~\ref{sec:sparsity-protocol}).

\subsubsection{Hàm mục tiêu và phương pháp tối ưu}
\label{sec:dfhf-objective}

Mỗi cặp $(u,i)$ trong tập giữ lại là một truy vấn huấn luyện. Ngữ cảnh của truy vấn là lịch sử bỏ đi chính sản phẩm dương (\emph{leave-one-out}):
\begin{equation}
q_{u\setminus i} = \frac{1}{(n_u-1)^{1/2}}\Bigl(\sum_{j\in H_u}\tilde S_j - \tilde S_i\Bigr).
\end{equation}
Hàm mất mát là cross-entropy của softmax trên toàn bộ danh mục, với các sản phẩm khác đã có trong lịch sử và các sản phẩm không có tương tác bị loại khỏi mẫu số, giống như cách đánh giá loại chúng khỏi danh sách gợi ý:
\begin{equation}
\ell(u,i) = -\log\frac{\exp\bigl(q_{u\setminus i}^\top\tilde V_i + b_i\bigr)}{\sum_{i'\in\mathcal{C}_u\cup\{i\}}\exp\bigl(q_{u\setminus i}^\top\tilde V_{i'} + b_{i'}\bigr)},
\label{eq:dfhf-loss}
\end{equation}
trong đó $\mathcal{C}_u$ là tập ứng viên trừ lịch sử của $u$. Dạng likelihood đa thức trên toàn danh mục này tương tự Mult-VAE~\cite{liang2018multvae}. Hàm mục tiêu cộng thêm phạt ridge có trọng số theo mức xuất hiện $\pi_i$ của sản phẩm trong các truy vấn ($\sum_i\pi_i=1$):
\begin{equation}
\mathcal{L} = \frac{1}{|E|}\sum_{(u,i)\in E}\ell(u,i) + \frac{\lambda}{2}\sum_i\pi_i\,\frac{\|V_i\|^2+\|S_i\|^2}{2} + \frac{\lambda_b}{2}\|b\|^2,
\end{equation}
với $\lambda=0{,}1$ và $\lambda_b=10^{-4}$. Sau mỗi bước cập nhật, mỗi hàng của $S$ và $V$ được chiếu về hình cầu bán kính 100; ràng buộc này hiếm khi có hiệu lực.

\subsubsection{Thuật toán huấn luyện và suy luận}
\label{sec:dfhf-algorithm}

DFHF được tối ưu bằng Adam với tốc độ học $3\cdot10^{-4}$, lô 2.048 truy vấn, trong đúng 200 epoch. Sau mỗi epoch, mô hình được đánh giá NDCG@20 trên validation và checkpoint có giá trị cao nhất được giữ lại. Đây là lựa chọn checkpoint chứ không phải dừng sớm: quá trình luôn chạy đủ 200 epoch. Khi huấn luyện, ngữ cảnh bỏ đi sản phẩm dương; khi suy luận, $q_u$ dùng toàn bộ lịch sử giữ lại theo \eqref{eq:dfhf-score}.

\begin{algorithm}[htbp]
\caption[Huấn luyện Diffusion-Filtered History Factorization]{Huấn luyện Diffusion-Filtered History Factorization.}
\label{alg:dfhf}
\KwIn{Tập cạnh giữ lại $E$; siêu tham số ở \tabref{tab:dfhf-hparams}}
\KwOut{Checkpoint $(S,V,b)$ có NDCG@20 validation cao nhất}
Khởi tạo $V\sim\mathcal{N}(0,1/64)$, $S\leftarrow V$ (sau đó cập nhật độc lập), $b\leftarrow0$\;
\For{epoch $=1,\dots,200$}{
  \ForEach{lô $\mathcal{Q}\subset E$ gồm 2.048 truy vấn}{
    Tính $\tilde S=FS$, $\tilde V=FV$ theo \eqref{eq:dfhf-filter}\;
    Tính ngữ cảnh $q_{u\setminus i}$ và mất mát \eqref{eq:dfhf-loss} cộng phạt\;
    Cập nhật Adam; chiếu các hàng của $S$, $V$ về bán kính 100\;
  }
  Đánh giá NDCG@20 trên validation; lưu checkpoint nếu tốt hơn\;
}
\end{algorithm}

\begin{table}[htbp]
\centering
\caption[Siêu tham số của Diffusion-Filtered History Factorization]{Siêu tham số của Diffusion-Filtered History Factorization. Cấu hình đã khóa của mô hình.}
\label{tab:dfhf-hparams}
\begin{tabular}{@{}llr@{}}
\toprule
Ký hiệu & Ý nghĩa & Giá trị \\
\midrule
--- & Số chiều embedding & 64 \\
--- & Tốc độ học Adam; kích thước lô & $3\cdot10^{-4}$; 2.048 \\
$\lambda$ & Phạt ridge cho embedding & 0,1 \\
$\lambda_b$ & Phạt cho vector chặn & $10^{-4}$ \\
--- & Bán kính chiếu & 100 \\
--- & Số mũ chuẩn hóa lịch sử & 0,5 \\
--- & Trọng số lịch sử đơn & 0 \\
--- & Nhiệt độ softmax & 1 \\
--- & Số epoch & 200 \\
\bottomrule
\end{tabular}
\end{table}

\subsection{Phân tích tính chất và chi phí tính toán}
\label{sec:complexity}

Mỗi mô hình được phân tích so với đại diện cùng nhóm, không so hai mô hình đề xuất với nhau (\tabref{tab:complexity}).

\textbf{GRNR so với EASE.} EASE~\cite{steck2019ease} cần nghịch đảo một ma trận $|\mathcal{I}|\times|\mathcal{I}|$ đầy đủ, tốn $O(|\mathcal{I}|^3)$ phép tính và $O(|\mathcal{I}|^2)$ bộ nhớ; với 44.843 sản phẩm, riêng ma trận float32 đã chiếm khoảng 7,5 GiB. GRNR thay bằng $|\mathcal{I}|$ hệ nhỏ kích thước tối đa $129\times129$, tốn $O(|\mathcal{I}|k^3)$ cho phần giải và lưu tối đa $k|\mathcal{I}|$ hệ số khác 0. Chi phí chính còn lại là ma trận đồng xuất hiện, cài đặt giới hạn tối đa $5\cdot10^7$ phần tử. Khi suy luận, cả hai đều là tích một vector lịch sử thưa với ma trận hệ số.

\textbf{DFHF so với LightGCN.} LightGCN~\cite{he2020lightgcn} học một embedding cho mỗi người dùng và mỗi sản phẩm, tức $64(|\mathcal{U}|+|\mathcal{I}|)$ tham số. DFHF chỉ học $2\cdot64|\mathcal{I}|+|\mathcal{I}|$ tham số, không phụ thuộc số người dùng. Đổi lại, mỗi epoch của DFHF tốn thêm chi phí cho softmax trên toàn bộ danh mục, $O(|E|\cdot|\mathcal{I}|\cdot 64)$, và hai lần áp toán tử $T$ trên mỗi bảng, $O(|E|\cdot 64)$ mỗi lần.

\begin{table}[htbp]
\centering
\caption[Tính chất của mỗi mô hình đề xuất so với đại diện cùng nhóm]{Tính chất của mỗi mô hình đề xuất so với đại diện cùng nhóm. Phân tích lý thuyết; số đo thực tế ở Mục~\ref{sec:grnr-results} và~\ref{sec:dfhf-results}.}
\label{tab:complexity}
\begin{tabularx}{\textwidth}{@{}lXXXX@{}}
\toprule
 & GRNR & EASE & DFHF & LightGCN \\
\midrule
Vector người dùng & Không & Không & Không & Có \\
Số tham số & $\le 129|\mathcal{I}|$ & $|\mathcal{I}|^2$ & $129|\mathcal{I}|$ & $64(|\mathcal{U}|+|\mathcal{I}|)$ \\
Cách học & $|\mathcal{I}|$ hệ Cholesky $\le129^2$ & Nghịch đảo $|\mathcal{I}|\times|\mathcal{I}|$ & Adam, 200 epoch & Adam, 200 epoch \\
Lịch sử mới & Chấm ngay & Chấm ngay & Chấm ngay & Cần huấn luyện lại \\
\bottomrule
\end{tabularx}
\end{table}
